\chapter{Hertesting en vergelijking van applicatieresultaten}
\label{ch:resultatenanalyse}

In dit hoofdstuk worden de oorspronkelijke experimenten herhaald onder zo vergelijkbaar mogelijke omstandigheden, om de impact van de aangebrachte wijzigingen te evalueren. Het hoofdstuk is zodanig gestructureerd dat de tests reproduceerbaar zijn en de vergelijking op een statistisch onderbouwde manier kan worden uitgevoerd.

\subsection{Hertesting via Lighthouse-audit}
\label{subsec:hertesting-lighthouse}

\subsubsection{Performantie}

De performantie van de applicatie werd geëvalueerd aan de hand van een Lighthouse-audit. Hierbij worden verschillende kernmetrics gebruikt die samen een beeld geven van de laadsnelheid, responsiviteit en visuele stabiliteit van de applicatie. De belangrijkste metrics zijn First Contentful Paint (FCP), Largest Contentful Paint (LCP), Total Blocking Time (TBT), Speed Index en Cumulative Layout Shift (CLS).

De gemeten waarden per pagina zijn weergegeven in onderstaande figuren en tabellen.

\begin{figure}[H]
\centering
\textbf{(a) Nulmeting}\\[1ex]
\begin{tikzpicture}    
  \begin{axis}[
    boxplot/draw direction=x,
    xlabel={First Contentful Paint (s)},
    xmin=5, xmax=7.5,
    xtick={5,5.5,6,6.5,7,7.5},
    ytick={1,2,3,4,5,6},
    yticklabels={Home, Nieuws, Campus, Lessenrooster, Deadline, Notificatie},
    ymajorgrids,
  ]
    \addplot[
      boxplot,
      fill=blue!20
    ] table[y index=0] {
      data
      7.2
      5.6
      5.6
      7.1
      7.1
      7.0
      7.0
      5.4
      5.6
      7.0
      5.6
      7.2
      5.7
      7.0
      5.5
      5.6
      5.3
      7.2
      7.0
      6.9
    };
    \addplot[
      boxplot,
      fill=blue!20
    ] table[y index=0] {
      data
      7.2
      7.0
      7.2
      7.1
      5.4
      7.2
      7.2
      7.2
      5.4
      7.0
      7.2
      5.4
      7.0
      7.2
      5.7
      5.8
      7.3
      6.9
      7.0
      7.2
    };
    \addplot[
      boxplot,
      fill=blue!20
    ] table[y index=0] {
      data
      5.3
      5.4
      6.8
      5.4
      5.4
      5.1
      5.4
      7.2
      5.2
      7.1
      7.0
      6.9
      7.1
      6.9
      7.1
    };
    \addplot[
      boxplot,
      fill=blue!20
    ] table[y index=0] {
      data
      5.4
      5.5
      7.2
      7.2
      7.2
      7.2
      5.4
      7.1
      7.2
      5.6
      5.5
      7.1
      6.9
      5.5
      5.5
    };
    \addplot[
      boxplot,
      fill=blue!20
    ] table[y index=0] {
      data
      6.9
      7.0
      7.2
      7.3
      5.2
      7.0
      7.0
      5.2
      5.2
      7.1
      7.1
      7.1
      7.1
      6.9
      7.0
    };
    \addplot[
      boxplot,
      fill=blue!20
    ] table[y index=0] {
      data
      5.5
      5.4
      5.3
      5.5
      7.2
      7.1
      5.4
      5.6
      7.1
      7.1
      5.5
      7.1
      6.9
      5.4
      5.5
      7.2
    };
  \end{axis}
\end{tikzpicture}

\vspace{1cm}

\textbf{(b) Hertesting}\\[1ex]
\begin{tikzpicture}
  \begin{axis}[
    boxplot/draw direction=x,
    xlabel={First Contentful Paint (s)},
    xmin=5, xmax=7.5,
    xtick={5,5.5,6,6.5,7,7.5},
    ytick={1,2,3,4,5,6},
    yticklabels={Home, Nieuws, Campus, Lessenrooster, Deadline, Notificatie},
    ymajorgrids,
  ]
    \addplot[
      boxplot,
      fill=red!20
    ] table[y index=0] {
      data
5.3
7.0
6.9
5.6
7.0
7.1
7.1
5.6
6.9
7.0
5.9
7.1
5.5
7.1
7.3
    };
    \addplot[
      boxplot,
      fill=red!20
    ] table[y index=0] {
      data
7.5
7.5
7.1
5.6
5.4
5.7
5.6
7.4
5.8
6.9
5.7
5.3
5.9
7.2
7.0
    };
    \addplot[
      boxplot,
      fill=red!20
    ] table[y index=0] {
      data
5.7
5.9
7.3
7.0
6.9
5.6
7.1
7.1
6.9
6.9
6.9
7.1
5.5
7.2
5.6
    };
    \addplot[
      boxplot,
      fill=red!20
    ] table[y index=0] {
      data
5.4
7.1
5.5
5.3
6.9
6.9
6.9
5.5
5.3
6.9
7.1
7.1
7.4
7.1
6.8
    };
    \addplot[
      boxplot,
      fill=red!20
    ] table[y index=0] {
      data
5.3
5.5
6.9
5.3
5.5
5.4
7.1
7.1
7.1
7.1
7.3
5.3
7.1
6.9
7.1
    };
    \addplot[
      boxplot,
      fill=red!20
    ] table[y index=0] {
      data
5.6
7.2
5.4
5.6
6.9
7.2
7.2
7.2
5.6
7.2
5.4
5.5
7.0
5.6
7.2
    };
  \end{axis}
\end{tikzpicture}
\caption{Vergelijking van First Contentful Paint (FCP): links de nulmeting, rechts de hertesting.}
\label{fig:fcp-comparison}
\end{figure}

De \textbf{First Contentful Paint (FCP)} blijven de resultaten in beide metingen relatief gelijkaardig. De mediane waarden bevinden zich op de meeste pagina’s tussen ongeveer 5,5 en 7 seconden. In de hertesting is echter zichtbaar dat verschillende boxplotten compacter worden, wat wijst op een stabielere rendering van de eerste zichtbare inhoud. Vooral bij de Home-, Nieuws- en Deadlinepagina is minder spreiding zichtbaar. Dit betekent dat de applicatie consistenter reageert na de optimalisaties.

\begin{figure}[H]
\centering
\textbf{(a) Nulmeting}\\[1ex]
\begin{tikzpicture}
	\begin{axis}[
		boxplot/draw direction=x,
		xlabel={Largest Contentful Paint (s)},
	  ytick={1,2},
		yticklabels={Home, Nieuws},
	  ymajorgrids,
		title={LCP Measurements}
	]
		\addplot[
			boxplot,
	  fill=blue!20
		] table[y index=0] {
			data
			90.9
			90.8
			91.0
			90.9
			90.9
			91.1
			90.8
			91.2
			90.9
			90.8
			90.9
			90.9
			91.0
			90.8
			90.9
		};
	\addplot[
			boxplot,
	  fill=blue!20
		] table[y index=0] {
			data
			92.1
			92.3
			92.3
			92.1
			92.3
			92.3
			92.4
			92.0
			92.3
			92.3
			92.3
			92.3
			92.2
			92.3
			92.6
		};
	\end{axis}
\end{tikzpicture}

\vspace{1cm}

\textbf{(b) Hertesting}\\[1ex]
\begin{tikzpicture}
  \begin{axis}[
    boxplot/draw direction=x,
    xlabel={Largest Contentful Paint (s)},
    xmin=0, xmax=100,
    xtick={0,10,20,30,40,50,60,70,80,90,100},
    ytick={1,2},
    yticklabels={Home, Nieuws},
    ymajorgrids,
  ]
    \addplot[boxplot, fill=red!20] table[y index=0] {
      data
      91.2
      91.1
      90.9
      90.9
      91.2
      7.5
      7.5
      6.0
      91.0
      7.7
      6.3
      90.8
      90.8
      7.7
      90.9
    };
    \addplot[boxplot, fill=red!20] table[y index=0] {
      data
      92.3
      90.0
      92.0
      92.8
      92.5
      92.6
      90.9
      91.7
      91.7
      92.1
      92.4
      92.3
      92.3
      92.1
      92.5
    };
  \end{axis}
\end{tikzpicture}
\caption{Vergelijking van Largest Contentful Paint (LCP) voor de Home- en Nieuwspagina: links de nulmeting, rechts de hertesting.}
\label{fig:lcp-comparison-home-nieuws}
\end{figure}

Voor de \textbf{Largest Contentful Paint (LCP)} blijven de resultaten tussen de nulmeting en de hertesting grotendeels stabiel. Bij de Nieuwspagina blijft zowel de mediaan als de spreiding vrijwel identiek, wat aangeeft dat de toegankelijkheidsimplementaties geen merkbare invloed hadden op het laden van de grootste zichtbare inhoud.

Bij de Homepagina is tijdens de hertesting een grotere spreiding zichtbaar. Dit wordt veroorzaakt door enkele metingen met aanzienlijk lagere waarden dan tijdens de nulmeting. Hierdoor ontstaat een bredere boxplot. Deze variatie wijst echter niet op een structurele verslechtering van de performantie. Integendeel, bepaalde metingen tonen zelfs een snellere rendering van de grootste content-elementen. De verschillen kunnen verklaard worden door externe factoren zoals caching of variaties in netwerkcondities tijdens de metingen.

\begin{figure}[H]
\centering
\textbf{(a) Nulmeting}\\[1ex]
\begin{tikzpicture}
	\begin{axis}[
		boxplot/draw direction=x,
		xlabel={Largest Contentful Paint (s)},
    xmin=5, xmax=8.5,
    xtick={5,5.5,6,6.5,7,7.5,8,8.5},
	  ytick={1,2,3,4},
		yticklabels={ Campus, Lessenrooster, Deadline, Notificatie},
	  ymajorgrids,
		title={LCP Measurements (continued) }
	]
	\addplot[
			boxplot,
	  fill=blue!20
		] table[y index=0] {
			data
			5.9
			6.0
			7.5
			6.0
			6.0
			5.9
			5.9
			7.5
			5.9
			7.5
			7.7
			7.5
			7.5
			7.6
			7.5
		};
	\addplot[
			boxplot,
	  fill=blue!20
		] table[y index=0] {
			data
			6.0
			5.9
			7.5
			7.7
			7.7
			7.5
			6.0
			7.5
			7.7
			6.0
			6.0
			7.5
			7.5
			6.0
			6.0
		};
	\addplot[
			boxplot,
	  fill=blue!20
		] table[y index=0] {
			data
			7.6
			7.7
			7.7
			7.7
			5.9
			7.7
			7.7
			5.9
			5.9
			7.5
			7.5
			7.5
			7.5
			7.5
			7.8
		};
	\addplot[
			boxplot,
	  fill=blue!20
		] table[y index=0] {
			data
			5.9
			5.9
			5.9
			6.0
			7.5
			7.5
			5.9
			6.0
			7.5
			7.5
			5.9
			7.5
			7.0
			5.9
			6.0
			7.5
		};
	\end{axis}
\end{tikzpicture}

\vspace{1cm}

\textbf{(b) Hertesting}\\[1ex]
\begin{tikzpicture}
  \begin{axis}[
    boxplot/draw direction=x,
    xlabel={Largest Contentful Paint (s)},
    xmin=5, xmax=8.5,
    xtick={5,5.5,6,6.5,7,7.5,8,8.5},
    ytick={1,2,3,4},
    yticklabels={Campus, Lessenrooster, Deadline, Notificatie},
    ymajorgrids,
  ]

    \addplot[boxplot, fill=red!20] table[y index=0] {
      data
      6.0
      6.3
      7.6
      7.8
      7.7
      6.3
      7.5
      7.5
      7.0
      7.5
      7.5
      7.5
      6.1
      7.7
      6.3
    };
    \addplot[boxplot, fill=red!20] table[y index=0] {
      data
      5.9
      7.5
      5.9
      5.9
      7.5
      7.5
      7.5
      5.9
      5.9
      7.5
      7.5
      7.5
      8.0
      7.5
      7.0
    };
    \addplot[boxplot, fill=red!20] table[y index=0] {
      data
      5.9
      5.9
      7.5
      5.9
      5.9
      5.9
      7.5
      7.5
      7.5
      7.5
      7.7
      5.9
      7.5
      7.5
      7.5
    };
    \addplot[boxplot, fill=red!20] table[y index=0] {
      data
      6.0
      7.8
      5.9
      6.0
      7.5
      7.7
      7.5
      7.7
      6.0
      7.7
      5.9
      5.9
      7.7
      6.0
      7.7
    };
  \end{axis}
\end{tikzpicture}
\caption{Vergelijking van Largest Contentful Paint (LCP) voor de overige pagina’s: links de nulmeting, rechts de hertesting.}
\label{fig:lcp-hertesting}
\end{figure}

Voor de overige pagina’s blijven de \textbf{Largest Contentful Paint (LCP)}-waarden eveneens grotendeels stabiel tussen de nulmeting en de hertesting. De mediane waarden bevinden zich bij beide meetmomenten voornamelijk tussen 6 en 8 seconden. Dit toont aan dat de toegankelijkheidsverbeteringen geen merkbare negatieve invloed hadden op het laden van de grootste zichtbare elementen binnen de applicatie.

Bij de Campus-, Lessenrooster- en Notificatiepagina blijven de boxplotten sterk vergelijkbaar qua spreiding en mediaan. Enkel kleine verschuivingen zijn zichtbaar, maar deze vallen binnen de normale variatie van performantiemetingen. De Deadlinepagina toont tijdens de hertesting zelfs een iets compactere spreiding, wat wijst op een stabielere rendering van content.

\begin{figure}[H]
\centering

\textbf{(a) Nulmeting}\\[1ex]

\begin{tikzpicture}
	\begin{axis}[
		boxplot/draw direction=x,
		xlabel={Speed Index (s)},
    xmin=7, xmax=18,
    xtick={7,8,9,10,11,12,13,14,15,16,17,18},
	ytick={1,2,3,4,5,6},
		yticklabels={Home, Nieuws, Campus, Lessenrooster, Deadline, Notificatie},
	ymajorgrids,
		title={Speed Index Measurements}
	]
		\addplot[
			boxplot,
	  fill=blue!20
		] table[y index=0] {
			data
			12.5
			9.0
			9.2
			9.2
			9.5
			12.1
			10.2
			9.9
			10.3
			10.2
			10.1
			9.5
			10.4
			10.1
			10.0
		};
	\addplot[
			boxplot,
	  fill=blue!20
		] table[y index=0] {
			data
			12.3
			10.5
			10.8
			11.0
			10.5
			10.6
			10.6
			10.5
			10.7
			10.3
			10.7
			11.1
			10.7
			11.0
			11.6
		};
	\addplot[
			boxplot,
	  fill=blue!20
		] table[y index=0] {
			data
			13.7
			10.6
			10.5
			10.5
			10.4
			11.1
			10.8
			10.9
			10.4
			10.1
			11.3
			11.1
			11.3
			10.8
			11.3
		};
	\addplot[
			boxplot,
	  fill=blue!20
		] table[y index=0] {
			data
			11.1
			10.1
			10.2
			9.6
			10.4
			10.4
			9.9
			9.9
			10.5
			9.6
			10.1
			10.1
			9.6
			8.9
			10.4
		};
	\addplot[
			boxplot,
	  fill=blue!20
		] table[y index=0] {
			data
			9.1
			8.1
			8.1
			8.2
			7.5
			8.3
			8.6
			8.3
			8.0
			9.0
			9.0
			8.4
			8.9
			9.0
			8.7
		};
	\addplot[
			boxplot,
	  fill=blue!20
		] table[y index=0] {
			data
			8.8
			7.8
			8.0
			8.4
			8.8
			8.1
			7.9
			8.0
			7.9
			8.4
			7.5
			7.9
			7.9
			7.6
			8.5
			8.8
		};
	\end{axis}
\end{tikzpicture}

\vspace{1cm}

\textbf{(b) Hertesting}\\[1ex]
\begin{tikzpicture}
  \begin{axis}[
    boxplot/draw direction=x,
    xlabel={Speed Index (s)},
    xmin=7, xmax=18,
    xtick={7,8,9,10,11,12,13,14,15,16,17,18},
    ytick={1,2,3,4,5,6},
    yticklabels={Home, Nieuws, Campus, Lessenrooster, Deadline, Notificatie},
    ymajorgrids,
  ]
    \addplot[boxplot, fill=red!20] table[y index=0] {
      data
      12.3
      9.0
      9.3
      9.4
      8.6
      9.0
      9.7
      8.5
      8.7
      9.0
      9.4
      9.2
      8.7
      9.0
      9.0
    };
    \addplot[boxplot, fill=red!20] table[y index=0] {
      data
      14.7
      7.9
      7.4
      8.5
      8.2
      7.7
      8.0
      8.8
      8.8
      7.9
      7.8
      8.3
      7.7
      8.0
      7.6
    };
    \addplot[boxplot, fill=red!20] table[y index=0] {
      data
      11.7
      12.6
      11.7
      11.9
      11.7
      11.9
      11.8
      11.8
      11.4
      11.6
      11.0
      11.1
      11.3
      11.8
      11.9
    };
    \addplot[boxplot, fill=red!20] table[y index=0] {
      data
      10.4
      10.0
      9.9
      9.6
      9.5
      9.8
      9.4
      10.4
      9.5
      9.7
      10.2
      11.0
      17.8
      10.8
      9.6
    };
    \addplot[boxplot, fill=red!20] table[y index=0] {
      data
      8.3
      8.2
      8.2
      8.1
      8.4
      8.0
      8.3
      8.3
      8.8
      8.7
      10.4
      8.1
      8.0
      8.6
      8.1
    };
    \addplot[boxplot, fill=red!20] table[y index=0] {
      data
      7.9
      8.3
      8.0
      8.0
      7.8
      8.7
      8.9
      8.4
      9.1
      8.2
      7.5
      8.7
      8.2
      9.1
      8.2
    };
  \end{axis}
\end{tikzpicture}
\caption{Vergelijking van Speed Index: links de nulmeting, rechts de hertesting.}
\label{fig:speed-index-comparison}
\end{figure}

De \textbf{Speed Index} toont een gemengd beeld. Op sommige pagina’s, zoals Home en Nieuws, is een lichte verbetering zichtbaar doordat de mediane waarde daalt en de boxplot compacter wordt. Dit betekent dat visuele inhoud sneller zichtbaar wordt voor de gebruiker. Bij andere pagina’s, zoals Campus, stijgt de mediane waarde licht tijdens de hertesting. Daarnaast verschijnt bij de Lessenroosterpagina een duidelijke uitschieter rond 17 seconden, wat wijst op een tijdelijke vertraging tijdens het renderen. Globaal kan gesteld worden dat de optimalisaties vooral een positieve invloed hadden op de consistentie van de visuele laadtijd, maar niet op alle pagina’s een directe snelheidswinst opleverden.

\begin{figure}[H]
\centering
\textbf{(a) Nulmeting}\\[1ex]
\begin{tikzpicture}
	\begin{axis}[
		boxplot/draw direction=x,
		xlabel={Total Blocking Time (ms)},
    xmin=0, xmax=18000,
    xtick={0,10000,12000,14000,16000,18000},
	  ytick={1,2,3,4,5,6},
		yticklabels={Home, Nieuws, Campus, Lessenrooster, Deadline, Notificatie},
	  ymajorgrids,
		title={TBT Measurements}
	]
		\addplot[
			boxplot,
	  fill=blue!20
		] table[y index=0] {
			data
			11740
			11980
			11170
			11300
			11150
			11720
			11420
			11400
			11670
			11310
			11650
			9240
			11410
			11240
			11320
		};
	\addplot[
			boxplot,
	  fill=blue!20
		] table[y index=0] {
			data
			9730
			11030
			11820
			10920
			11100
			10770
			10750
			11340
			11710
			10630
			10400
			10930
			11110
			11620
			15560
		};
	\addplot[
			boxplot,
	  fill=blue!20
		] table[y index=0] {
			data
			10150
			10590
			11030
			11350
			10730
			10620
			10650
			10810
			11410
			10150
			10420
			10820
			10540
			11490
			9980
		};
	\addplot[
			boxplot,
	  fill=blue!20
		] table[y index=0] {
			data
			10000
			9760
			9560
			9710
			9680
			9660
			10380
			10800
			10200
			10320
			9930
			9230
			10100
			9510
			9930
		};
	\addplot[
			boxplot,
	  fill=blue!20
		] table[y index=0] {
			data
			11740
			9730
			10030
			10180
			9040
			9080
			9750
			9520
			10240
			9920
			10010
			9440
			10040
			10000
			11000
		};
	\addplot[
			boxplot,
	  fill=blue!20
		] table[y index=0] {
			data
			11420
			9510
			9040
			9850
			10640
			9580
			10680
			10040
			10540
			10210
			9770
			10050
			10850
			11000
			10450
			9660
		};
	\end{axis}
\end{tikzpicture}

\vspace{1cm}

\textbf{(b) Hertesting}\\[1ex]
\begin{tikzpicture}
  \begin{axis}[
    boxplot/draw direction=x,
    xlabel={Total Blocking Time (ms)},
    xmin=0, xmax=18000,
    xtick={0,10000,12000,14000,16000,18000},
    ytick={1,2,3,4,5,6},
    yticklabels={Home, Nieuws, Campus, Lessenrooster, Deadline, Notificatie},
    ymajorgrids,
  ]
    \addplot[boxplot, fill=red!20] table[y index=0] {
      data
      12040
      11010
      11660
      12990
      10740
      10910
      10470
      9830
      10790
      11640
      10800
      11670
      10860
      10690
      10740
    };
    \addplot[boxplot, fill=red!20] table[y index=0] {
      data
      11500
      11420
      10520
      10900
      11340
      10710
      10470
      10780
      11110
      10620
      10730
      11170
      10290
      11420
      11240
    };
    \addplot[boxplot, fill=red!20] table[y index=0] {
      data
      9650
      12290
      11010
      10770
      10360
      10250
      12020
      11240
      11270
      11890
      11120
      11260
      10460
      11240
      10250
    };
    \addplot[boxplot, fill=red!20] table[y index=0] {
      data
      10810
      10650
      10120
      9400
      8860
      10010
      9290
      9670
      9980
      9600
      8940
      10770
      15640
      10400
      10710
    };
    \addplot[boxplot, fill=red!20] table[y index=0] {
      data
      12010
      11720
      10560
      11390
      12460
      8940
      10550
      9980
      9910
      10660
      17280
      9610
      10010
      11470
      10660
    };
    \addplot[boxplot, fill=red!20] table[y index=0] {
      data
      9150
      10330
      9410
      11090
      10390
      10400
      10630
      11080
      12520
      10450
      10180
      11060
      9300
      12520
      10450
    };
  \end{axis}
\end{tikzpicture}
\caption{Vergelijking van Total Blocking Time (TBT): links de nulmeting, rechts de hertesting.}
\label{fig:tbt-comparison}
\end{figure}

Voor de \textbf{Total Blocking Time (TBT)} blijven de resultaten relatief vergelijkbaar tussen beide metingen. De mediane waarden blijven voor de meeste pagina’s rond 9,000 tot 12,000 milliseconden liggen. Wel valt op dat sommige pagina’s in de hertesting een iets compactere spreiding vertonen, wat duidt op een stabielere verwerking. Tegelijk blijven enkele sterke uitschieters aanwezig, zoals bij de Deadlinepagina en Lessenroosterpagina, waar pieken tot boven 15,000 milliseconden voorkomen. Hierdoor kan geconcludeerd worden dat de optimalisaties slechts een beperkte invloed hadden op het verminderen van blocking time.

\begin{table}[H]
\centering
\caption[CLS-waarden per pagina: nulmeting en hertesting]{Overzicht van de Cumulative Layout Shift-waarden per pagina tijdens de nulmeting en hertesting.}
\label{tab:cls-hertesting-per-pagina}
\begin{tabular}{lc}
\hline
\textbf{Pagina} & \textbf{CLS (nulmeting / hertesting)} \\
\hline
Homepagina & 0.079 / 0.079 \\
Nieuwspagina & 0.078 / 0.078 \\
Campuspagina & 0.080 / 0.080 \\
Lessenroosterpagina & 0.079 / 0.079 \\
Deadlinepagina & 0.080 / 0.080 \\
Notificatiepagina & 0.079 / 0.079 \\
\hline
\end{tabular}
\end{table}

De resultaten van de \textbf{Cumulative Layout Shift (CLS)} tonen daarentegen een zeer stabiel beeld. Alle pagina’s behalen waarden rond 0,078 tot 0,080, wat ruim binnen de aanbevolen grens van 0,1 ligt. Dit betekent dat de lay-out van de applicatie visueel stabiel blijft tijdens het laden en dat gebruikers weinig onverwachte verschuivingen ervaren. De optimalisaties hadden hier weinig zichtbare impact omdat de applicatie reeds goede CLS-scores behaalde tijdens de nulmeting.


\paragraph{Algemene conclusie}

Globaal tonen de resultaten aan dat de performantie van de applicatie na de toegankelijkheidsoptimalisaties grotendeels stabiel is gebleven. Voor de meeste metrics, waaronder FCP, LCP en TBT, worden geen significante verbeteringen of degradaties vastgesteld.

Enkele beperkte verbeteringen in spreiding en consistentie zijn zichtbaar, voornamelijk bij FCP en in beperkte mate bij Speed Index. Deze effecten zijn echter niet uniform over alle pagina’s en metrics heen.

De CLS-waarden blijven volledig stabiel en binnen de aanbevolen grenswaarden. De resultaten suggereren bijgevolg dat de uitgevoerde toegankelijkheidsaanpassingen geen negatieve impact hebben gehad op de performantie van de applicatie.

\subsubsection{toegankelijkheid}

De toegankelijkheid van de applicatie werd geëvalueerd via een Lighthouse-audit. Deze metric geeft een indicatie van hoe goed de applicatie bruikbaar is voor gebruikers met verschillende beperkingen, zoals visuele of motorische beperkingen. De score wordt uitgedrukt op een schaal van 100, waarbij een hogere score duidt op betere toegankelijkheid.

De gemeten waarden per pagina zijn weergegeven in onderstaande tabel.

\begin{table}[h]
\centering
\caption[Toegankelijkheidsscore per pagina]{Overzicht van de Lighthouse-toegankelijkheidsscore per pagina.}
\label{tab:toegankelijkheid-score-per-pagina}
\begin{tabular}{lcc}
\hline
\textbf{Pagina} & \textbf{Totale score (/100)} & \textbf{Zone} \\
\hline
Homepagina & 96 & Goed \\
Nieuwspagina & 100 & Goed \\
Campuspagina & 100 & Goed \\
Lessenroosterpagina & 100 & Goed \\
Deadlinepagina & 100 & Goed \\
Notificatiepagina & 100 & Goed \\
\hline
\end{tabular}
\end{table}

Uit de resultaten blijkt dat de applicatie in het algemeen zeer goed scoort op toegankelijkheid. Vijf van de zes pagina's behalen een perfecte score van 100/100, wat wijst op een correcte implementatie van semantische HTML, voldoende structuur en een goede ondersteuning voor assistieve technologieën.

Enkel op de homepagina werd een lichte afwijking vastgesteld met een score van 96/100. Lighthouse detecteerde hier een klein probleem met onvoldoende contrast tussen bepaalde tekst- en achtergrondkleuren. Dit kan de leesbaarheid voor gebruikers met visuele beperkingen licht beïnvloeden, maar heeft geen kritieke impact op de algemene bruikbaarheid van de applicatie.

Hoewel de impact beperkt is, wordt dit beschouwd als een verbeterpunt binnen de verdere optimalisatie van de gebruikersinterface.

\subsection{Manuele test: nulmeting vs hertesting (toegankelijkheid)}
In aanvulling op de geautomatiseerde Lighthouse-audit werd de toegankelijkheid ook manueel geëvalueerd, met specifieke aandacht voor toetsenbordnavigatie en ondersteuning door schermlezers. Deze sectie bespreekt de verschillen tussen de eerder uitgevoerde nulmeting en de hertesting.

Tijdens de nulmeting werd vastgesteld dat de applicatie aanzienlijke tekortkomingen vertoonde op het vlak van toetsenbordtoegankelijkheid. De navigatie via de tabtoets was beperkt en inconsistent, waardoor gebruikers niet op een logische manier doorheen de interface konden navigeren. Daarnaast werden verschillende interactieve elementen niet correct geïnterpreteerd door de schermlezer NVDA of ontbraken duidelijke aankondigingen volledig. Dit bemoeilijkte zowel de toegang tot als het begrip van de inhoud.

De resultaten van de hertesting tonen aan dat deze problemen in belangrijke mate zijn verholpen. De focusvolgorde werd herwerkt zodat deze overeenkomt met de semantische en visuele structuur van de applicatie. Hierdoor is een consistente en voorspelbare navigatie via de tabtoets mogelijk geworden.

Verder worden interface-elementen nu correct herkend en aangekondigd door NVDA. Dit omvat onder andere het voorlezen van labels, rollen en statussen van interactieve componenten, waardoor gebruikers een beter inzicht krijgen in de functionaliteit en structuur van de applicatie.

De vergelijking tussen beide meetmomenten wijst op een duidelijke verbetering in de praktische toegankelijkheid van de applicatie. Waar de nulmeting wees op structurele problemen in navigatie en interpretatie, bevestigt de hertesting dat de doorgevoerde aanpassingen hebben geleid tot een aanzienlijk betere ondersteuning van assistieve technologieën en toetsenbordgebruik.
